\documentclass[letterpaper,12pt]{report}
\usepackage{thesis}

\begin{document}

% ----- Páginas preliminares en números romanos -----
\pagenumbering{roman}

% ----- Portada -----
\thesistitle
{Proyecto de Investigación: \\
	Análisis, Diseño y Construcción de un Sistema Operativo desde Cero}
{\emph{\copyright\ Magaña Osorio Jhoel Fabrizzio}\\
	\emph{\copyright\ Colque Quispe Fidel Enrique}\\
	\emph{\copyright\ Montalvo Sol\'orzano Rosy Aurely}\\
	\emph{\copyright\ Quispe Llavilla Jhon Andherson}}
{Docente: \emph{Ugarte Rojas Hector Ugarte}}
{Universidad Nacional de San Antonio Abad del Cusco\\
	Facultad de Ingeniería Eléctrica, Electrónica, Informática y Mecánica \\
	Escuela Profesional de Ingeniería Informática y de Sistemas}
{\emph{Cusco - Perú\\2025}}

% ----- Índices (TOC/LOT/LOF) -----
% Asegúrate de que contents.tex, tables.tex y figures.tex
% contengan \tableofcontents, \listoftables y \listoffigures respectivamente.
\include{contents}
\include{tables}
\include{figures}

% ----- Cuerpo en arábigos -----
\pagenumbering{arabic}

% =============================
% E1: Investigación y documentación
% =============================
\input{sections/e1-01-introduccion}
\input{sections/e1-02-marco-teorico}
\input{sections/e1-03-revision-sistemas}
\input{sections/e1-04-comparacion-tecnica}

% =============================
% E2: Propuestas y comparación
% =============================
\input{sections/e2-01-propuestas-analizadas}
\input{sections/e2-02-comparacion-tecnica}
\input{sections/e2-03-analisis-critico}

% =============================
% E3: Selección, entorno y evidencia
% =============================
\input{sections/e3-01-propuesta-seleccionada}
\input{sections/e3-02-argumentos-tecnicos}
\input{sections/e3-03-argumentos-pedagogicos}
\input{sections/e3-04-diseno-sistema}
\input{sections/e3-05-herramientas-entorno}
\section{Conclusiones}

El análisis desarrollado permitió comprender la diversidad de enfoques existentes en el diseño e implementación de sistemas operativos, evidenciando que no existe una única arquitectura óptima, sino soluciones adaptadas a distintos propósitos, contextos y niveles de complejidad. La comparación entre sistemas monolíticos, microkernel, híbridos y experimentales facilitó la identificación de las ventajas y limitaciones de cada propuesta, así como su impacto en aspectos clave como el rendimiento, la seguridad, la estabilidad y la mantenibilidad.

Asimismo, el estudio de sistemas operativos consolidados y alternativos puso de manifiesto la importancia de la documentación, la comunidad de desarrollo y las herramientas disponibles como factores determinantes para su adopción en entornos académicos. Los sistemas con documentación clara, código bien estructurado y soporte activo ofrecen mayores oportunidades de aprendizaje, permitiendo a los estudiantes analizar internamente componentes críticos como la gestión de procesos, memoria, sistemas de archivos y mecanismos de comunicación.

La evaluación comparativa realizada permitió justificar la selección de SerenityOS como sistema de referencia para el desarrollo del trabajo, debido a su equilibrio entre complejidad técnica y valor educativo. Su arquitectura bien definida, el uso predominante de C++ moderno, la modularidad del código y la disponibilidad de herramientas de compilación, ejecución y depuración lo convierten en una plataforma adecuada para el estudio práctico de los conceptos fundamentales de sistemas operativos.

Finalmente, este trabajo establece una base sólida para futuras etapas de análisis y desarrollo, orientadas a la experimentación, extensión o implementación de nuevos módulos sobre el sistema seleccionado. La metodología aplicada y los criterios definidos pueden reutilizarse en investigaciones posteriores, contribuyendo a una comprensión más profunda del funcionamiento interno de los sistemas operativos y fortaleciendo la formación académica en esta área fundamental de la ingeniería informática.


% ----- Bibliografía (BibTeX + apacite) -----
\input{bib}

% ----- Anexos / Apéndices -----
\appendix
% ==================================================
% Anexos / Apéndices (Entregable 3)
% ==================================================
\chapter{Anexos}\label{chap:anexos}

En esta sección se incluyen materiales complementarios que respaldan el trabajo realizado en el Entregable 3, con el objetivo de evidenciar el entorno utilizado, el proceso de ejecución/emulación y los elementos técnicos más relevantes.

\section{Diagramas técnicos}
Esta subsección presenta diagramas que sintetizan el flujo de construcción, instalación y ejecución/emulación de SerenityOS, así como su organización general a nivel de componentes.

\begin{figure}[H]
\centering
\IfFileExists{figures/anexos/diagrama-flujo-build-run.png}{%
\includegraphics[width=0.95\textwidth]{figures/anexos/diagrama-flujo-build-run.png}%
}{%
\fbox{\parbox{0.9\textwidth}{\centering\vspace{2.5cm}
	extbf{Imagen pendiente:} \texttt{figures/anexos/diagrama-flujo-build-run.png}\\
\vspace{2.5cm}}}%
}
\caption{Diagrama del flujo general de trabajo para construir y ejecutar SerenityOS en el host: instalación de dependencias, construcción de toolchain, generación de artefactos y arranque del sistema en QEMU.}\label{fig:anexo-flujo-build-run}
\end{figure}

\begin{figure}[H]
\centering
\IfFileExists{figures/anexos/diagrama-componentes-serenityos.png}{%
\includegraphics[width=0.95\textwidth]{figures/anexos/diagrama-componentes-serenityos.png}%
}{%
\fbox{\parbox{0.9\textwidth}{\centering\vspace{2.5cm}
	extbf{Imagen pendiente:} \texttt{figures/anexos/diagrama-componentes-serenityos.png}\\
\vspace{2.5cm}}}%
}
\caption{Diagrama de componentes a alto nivel: Kernel, bibliotecas base, servicios y aplicaciones de usuario, con sus relaciones principales.}\label{fig:anexo-componentes-serenityos}
\end{figure}

\begin{figure}[H]
\centering
\IfFileExists{figures/anexos/diagrama-estructura-repo.png}{%
\includegraphics[width=0.95\textwidth]{figures/anexos/diagrama-estructura-repo.png}%
}{%
\fbox{\parbox{0.9\textwidth}{\centering\vspace{2.5cm}
	extbf{Imagen pendiente:} \texttt{figures/anexos/diagrama-estructura-repo.png}\\
\vspace{2.5cm}}}%
}
\caption{Diagrama de la estructura general del repositorio de SerenityOS (directorios principales y propósito de cada uno) como apoyo a la navegación y documentación técnica.}\label{fig:anexo-estructura-repo}
\end{figure}

\section{Capturas del entorno de desarrollo}
Se incluyen capturas representativas del entorno utilizado para editar el documento, compilar el sistema y ejecutar la emulación.

\begin{figure}[H]
\centering
\IfFileExists{figures/anexos/captura-vscode-workspace.png}{%
\includegraphics[width=0.95\textwidth]{figures/anexos/captura-vscode-workspace.png}%
}{%
\fbox{\parbox{0.9\textwidth}{\centering\vspace{2.5cm}
	extbf{Imagen pendiente:} \texttt{figures/anexos/captura-vscode-workspace.png}\\
\vspace{2.5cm}}}%
}
\caption{Captura del entorno de edición (Visual Studio Code) con la estructura del repositorio y los archivos de configuración relevantes para el desarrollo y la documentación.}\label{fig:anexo-captura-vscode}
\end{figure}

\begin{figure}[H]
\centering
\IfFileExists{figures/anexos/captura-build-run-terminal.png}{%
\includegraphics[width=0.95\textwidth]{figures/anexos/captura-build-run-terminal.png}%
}{%
\fbox{\parbox{0.9\textwidth}{\centering\vspace{2.5cm}
	extbf{Imagen pendiente:} \texttt{figures/anexos/captura-build-run-terminal.png}\\
\vspace{2.5cm}}}%
}
\caption{Captura de terminal durante el flujo de compilación y ejecución mediante el script \texttt{Meta/serenity.sh}.}\label{fig:anexo-captura-terminal}
\end{figure}

\begin{figure}[H]
\centering
\IfFileExists{figures/anexos/captura-qemu-serenityos.png}{%
\includegraphics[width=0.95\textwidth]{figures/anexos/captura-qemu-serenityos.png}%
}{%
\fbox{\parbox{0.9\textwidth}{\centering\vspace{2.5cm}
	extbf{Imagen pendiente:} \texttt{figures/anexos/captura-qemu-serenityos.png}\\
\vspace{2.5cm}}}%
}
\caption{Captura de la ventana de QEMU con SerenityOS en ejecución, evidenciando el arranque y el entorno gráfico del sistema.}\label{fig:anexo-captura-qemu}
\end{figure}

\section{Fragmentos de código ilustrativo}
Los siguientes fragmentos se presentan como referencia práctica. Se basan en comandos habituales del ecosistema GNU/Linux (consultables mediante \texttt{man} o \texttt{--help}) y en el flujo recomendado por SerenityOS para construir y ejecutar el sistema en modo emulado.

\subsection{Instalación de dependencias (Arch Linux)}
Ejemplo de instalación de dependencias mínimas en Arch Linux (x86\_64):

\begin{verbatim}
sudo pacman -S --needed \
	base-devel cmake curl mpfr libmpc gmp e2fsprogs ninja qemu-desktop \
	qemu-system-aarch64 ccache rsync unzip

# Opcional:
# - fuse2fs (imagenes sin root)
# - clang llvm llvm-libs (para compilar con Clang)
\end{verbatim}

\subsection{Clonado del repositorio y ejecución del flujo principal}
Comandos de referencia para obtener el código, compilar y ejecutar mediante emulación (QEMU).

\begin{verbatim}
git clone https://github.com/SerenityOS/serenity.git
cd serenity

# Compila y ejecuta SerenityOS en QEMU.
Meta/serenity.sh run

# Compila sin arrancar la VM.
Meta/serenity.sh build
\end{verbatim}

\subsection{Fragmento ilustrativo de opciones de QEMU}
En sistemas donde se ejecuta QEMU directamente (en lugar de invocar el wrapper de SerenityOS), la sintaxis general es:

\begin{verbatim}
qemu-system-x86_64 [opciones]
\end{verbatim}

Ejemplo ilustrativo (parámetros frecuentes en entornos de emulación):

\begin{verbatim}
qemu-system-x86_64 \
	-m 2048 \
	-smp 4 \
	-display gtk \
	-serial stdio \
	-drive file=DISK_IMAGE.img,format=raw,if=virtio
\end{verbatim}

\section{Manual de usuario básico e instrucciones de ejecución}
Este manual resume los pasos para reproducir el entorno y ejecutar SerenityOS en modo emulado, así como acciones básicas de uso durante la sesión.

\subsection{Requisitos previos}
\begin{itemize}
\item Sistema host GNU/Linux de 64 bits (se utilizó Arch Linux x86\_64).
\item Conectividad a Internet (descargas iniciales del flujo de build).
\item Paquetes de build y QEMU instalados (ver fragmento de dependencias).
\end{itemize}

\subsection{Ejecución paso a paso}
\begin{enumerate}
\item Clonar el repositorio de SerenityOS y entrar al directorio de trabajo.
\item Ejecutar \texttt{Meta/serenity.sh run}. En la primera ejecución es esperable mayor tiempo por descarga de bases y construcción de la toolchain.
\item Verificar que se abra la ventana de QEMU y que el sistema complete el arranque.
\end{enumerate}

\subsection{Interacción básica durante la emulación}
\begin{itemize}
\item La ventana de QEMU captura teclado y ratón al interactuar con el sistema invitado.
\item Para liberar el cursor, QEMU utiliza una combinación de teclas conocida como \emph{grab/ungrab} (la combinación específica depende del backend de QEMU y del sistema; puede consultarse en el menú de QEMU o en la documentación del front-end utilizado).
\item Para finalizar una sesión, puede cerrarse la ventana de QEMU o detenerse el proceso desde la terminal (según cómo se haya lanzado la ejecución).
\end{itemize}

\subsection{Evidencias esperadas}
Como evidencia mínima de ejecución, se recomienda adjuntar:
\begin{itemize}
\item Captura del comando \texttt{Meta/serenity.sh run} finalizando la compilación e iniciando la ejecución emulada.
\item Captura de SerenityOS ya iniciado (pantalla principal o una aplicación abierta).
\end{itemize}


\end{document}
